% =============================================================
% บทที่ 1 - บทนำ
% =============================================================

% -----------------------------------------------------------
\section{ความเป็นมาและความสำคัญของปัญหา}
% -----------------------------------------------------------

ในปัจจุบัน ข่าวสารเป็นข้อมูลสำคัญที่ช่วยให้ประชาชนรับรู้และติดตามเหตุการณ์ได้อย่างทันท่วงที โดยเฉพาะอย่างยิ่งในยุคที่ข้อมูลมีการเปลี่ยนแปลงตลอดเวลา อย่างไรก็ตาม ผู้อ่านที่ต้องการติดตามข่าวสารจากหลายแหล่งมักประสบปัญหาหลายประการ ได้แก่ ต้องเปิดเว็บไซต์หรือแอปพลิเคชันหลายแห่งเพื่อเปรียบเทียบข่าวเดียวกัน ทำให้เสียเวลาและเกิดความยุ่งยากในการติดตามสถานการณ์ที่เปลี่ยนแปลงอย่างรวดเร็ว นอกจากนี้ ข่าวแต่ละฉบับยังมีเนื้อหายาวและมีข้อมูลจำนวนมาก ทำให้ผู้อ่านต้องใช้เวลาในการอ่านเพื่อสกัดประเด็นสำคัญ รวมถึงยังต้องอ่านข่าวจากแหล่งภาษาอังกฤษเพื่อให้ได้มุมมองที่หลากหลาย อีกทั้งข่าวสารที่มีปริมาณมากและหลากหลายหมวดหมู่ยังทำให้การคัดกรองและวิเคราะห์ข่าวเป็นเรื่องท้าทาย

การนำเทคโนโลยีเว็บสแครปปิง (Web Scraping) และ RSS Feed มาใช้ในการดึงข้อมูลข่าวจากเว็บไซต์ข่าวอัตโนมัติ ตลอดจนการใช้เทคนิคการประมวลผลภาษาธรรมชาติ (Natural Language Processing: NLP) ในการจำแนกหมวดหมู่ข่าวตามประเภท และการใช้โมเดลภาษาขนาดใหญ่ (Large Language Model: LLM) ในการสรุปเนื้อหา จะช่วยลดภาระของผู้อ่านในการเปิดหลายเว็บไซต์และอ่านบทความฉบับเต็ม อีกทั้งยังสามารถนำเสนอข้อมูลสำคัญที่คัดสรรและวิเคราะห์แล้วแบบเรียลไทม์ได้ในที่เดียว

จากปัญหาดังกล่าว โครงงานนี้จึงมีแนวคิดในการพัฒนาระบบจำแนกและวิเคราะห์ข่าวสารอัตโนมัติด้วยเทคนิคการประมวลผลภาษาธรรมชาติแบบเฉพาะประเภท (Category-Specific NLP Analytics Pipeline) โดยระบบจะทำการดึงข่าวจากแหล่งข่าวเป็นระยะแบบอัตโนมัติ จำแนกหมวดหมู่ข่าวด้วยเทคนิค NLP และโมเดลแมชชีนเลิร์นนิง สรุปเนื้อหาด้วยโมเดลภาษาขนาดใหญ่ และแสดงผลข่าวพร้อมการวิเคราะห์บนเว็บแอปพลิเคชันแบบเรียลไทม์ผ่าน WebSocket เพื่อให้ผู้ใช้สามารถติดตามข่าวสารได้อย่างสะดวก รวดเร็ว และเหมาะสมกับความสนใจของแต่ละบุคคล

% -----------------------------------------------------------
\section{วัตถุประสงค์}
% -----------------------------------------------------------

\begin{enumerate}[label=\arabic*)]
    \item เพื่อศึกษาและวิเคราะห์เทคนิคการดึงข้อมูลข่าวจากเว็บไซต์ข่าวด้วยวิธีเว็บสแครปปิงและการอ่าน RSS Feed รวมถึงเทคนิคการประมวลผลภาษาธรรมชาติสำหรับการจำแนกหมวดหมู่ข่าวตามประเภท
    \item เพื่อออกแบบและพัฒนาระบบจำแนกและวิเคราะห์ข่าวสารอัตโนมัติด้วยเทคนิคการประมวลผลภาษาธรรมชาติแบบเฉพาะประเภท (Category-Specific NLP Analytics Pipeline) ครอบคลุมการดึงข่าว การจำแนกหมวดหมู่ การสรุปเนื้อหาด้วยโมเดลภาษาขนาดใหญ่ อัลกอริทึมแนะนำข่าว และการแสดงผลแบบเรียลไทม์
    \item เพื่อทดสอบและประเมินประสิทธิภาพของระบบในด้านความถูกต้องของการจำแนกหมวดหมู่ข่าว ความถูกต้องของการสรุปข่าว ความแม่นยำของอัลกอริทึมแนะนำข่าว และความสามารถในการแสดงผลข่าวแบบเรียลไทม์
\end{enumerate}

% -----------------------------------------------------------
\section{ขอบเขตและข้อจำกัดของระบบ}
% -----------------------------------------------------------

\subsection{ขอบเขต}

ระบบเป็นเว็บแอปพลิเคชันแบบเต็มรูปแบบ (Full-stack) ที่มีขอบเขตการทำงาน ดังนี้

\begin{enumerate}[label=\arabic*)]
    \item \textbf{การรวบรวมข่าวอัตโนมัติ:} ระบบสามารถดึงข่าวจากแหล่งข่าวออนไลน์หลายแห่ง โดยใช้เทคนิคเว็บสแครปปิง การอ่าน RSS Feed และการจำลองเบราว์เซอร์ด้วย Playwright สำหรับเว็บไซต์ที่ใช้ JavaScript
    \item \textbf{การสรุปข่าวและวิเคราะห์อารมณ์:} ระบบส่งเนื้อหาข่าวในรูปแบบ Markdown ให้โมเดลภาษาขนาดใหญ่ผ่าน API ของมหาวิทยาลัยขอนแก่น (KKU Gen.AI) เพื่อสร้างบทสรุป ประเด็นสำคัญ อารมณ์ของข่าว (Sentiment) และคำสำคัญ โดยสรุปเป็นภาษาเดียวกับต้นฉบับข่าว
    \item \textbf{การจำแนกหมวดหมู่ข่าว 9 หมวดหมู่หลัก:} ระบบจำแนกข่าวออกเป็น 9 หมวดหมู่หลัก ได้แก่ การเมือง (politics), เศรษฐกิจ (economy), เทคโนโลยี (tech), สุขภาพ (health), สิ่งแวดล้อม (environment), กีฬา (sports), บันเทิง (entertainment), สังคม (society) และต่างประเทศ (world) โดยใช้การตัดคำภาษาไทยด้วย PyThaiNLP ร่วมกับโมเดลการเรียนรู้ของเครื่องแบบหลายชั้น (3-Tier Cascade Classifier: Rule Cues, Calibrated LinearSVC และ WangchanBERTa)
    \item \textbf{การจัดเก็บและวงจรชีวิตของข้อมูล (Data Lifecycle \& Retention):} ข้อมูลข่าวและผลการวิเคราะห์ทั้งหมดจะถูกจัดเก็บรวบรวมอย่างต่อเนื่องนับตั้งแต่วันที่ Instance Server เริ่มทำงาน (Instance Startup Date) โดยจัดเก็บไว้ภายใน Docker Container และ Local Storage บนเครื่องแม่ข่าย AWS EC2
    \item \textbf{การจัดกลุ่มประเด็นข่าวหลายสำนัก (Multi-Source Story Clustering):} ระบบสามารถตรวจจับ เปรียบเทียบความคล้ายคลึง และจัดกลุ่มบทความข่าวจากหลากหลายสำนักข่าวที่มีเนื้อหาเกี่ยวกับเหตุการณ์เดียวกันภายในกรอบเวลาไม่เกิน 36 ชั่วโมงเข้าเป็นกลุ่มประเด็นเดียวกัน (Story Cluster) เพื่อป้องกันการแสดงข่าวซ้ำซ้อน
    \item \textbf{การคัดเลือกข่าวตัวแทนและการหาฉันทามติหมวดหมู่ระดับกลุ่ม (Cluster Representative \& Category Consensus):} ในแต่ละกลุ่มข่าว ระบบจะคัดเลือกบทความที่มีคะแนนความนิยม ($\text{TrendingScore}$) สูงสุดขึ้นมาเป็นตัวแทนบทความหลักในการแสดงผล พร้อมแสดงตัวบ่งชี้จำนวนสำนักข่าวที่รายงานร่วมกัน และกำหนดหมวดหมู่ของกลุ่มข่าวด้วยหลักการมติเสียงข้างมาก (Majority Voting) เพื่อแก้ไขปัญหาความไม่สอดคล้องของหมวดหมู่ข่าวข้ามสำนักพิมพ์
    \item \textbf{ขอบเขตความเป็นส่วนตัวและการไหลของข้อมูลผู้ใช้ (User Privacy \& Data Flow):} ระบบบันทึกเฉพาะรหัสเซสชันแบบไม่ระบุตัวตน (Anonymous Cookie ID) และสถิติพฤติกรรมการอ่านข่าวแบบภาพรวม (Non-PII Aggregate Engagement: การคลิกเปิดอ่าน, การขยายอ่านบทสรุป, การบุ๊กมาร์ก) บนเบราว์เซอร์ของผู้ใช้เพื่อส่งมาคำนวณคะแนนความนิยม โดยไม่มีการจัดเก็บข้อมูลระบุตัวตนส่วนบุคคล (PII) เช่น ชื่อ-นามสกุล อีเมล หรือหมายเลข IP และไม่มีการส่งต่อข้อมูลพฤติกรรมผู้ใช้ไปยังโมเดลภาษาขนาดใหญ่ภายนอกหรือบุคคลภายนอกใด ๆ
    \item \textbf{การแสดงผลแบบเรียลไทม์:} ระบบแสดงรายการข่าวบนเว็บแอปพลิเคชันพร้อมสรุปเนื้อหา ภาพประกอบ และหมวดหมู่ และแจ้งเตือนข่าวใหม่ทันทีผ่าน WebSocket (Socket.IO) โดยไม่ต้องรีเฟรชหน้าเว็บ
    \item \textbf{การอ่านข่าวในตัวเว็บ (In-app Reader Modal):} ระบบมีหน้าต่างแสดงผลบทความข่าวและบทวิเคราะห์ในตัวเว็บแอปพลิเคชัน เพื่อให้ผู้ใช้อ่านข่าวได้อย่างต่อเนื่องโดยไม่ต้องเปิดแท็บใหม่
    \item \textbf{การจัดเก็บข้อมูลและสถาปัตยกรรมแบบแยกชั้น:} ระบบจัดเก็บข้อมูลข่าวและสถิติการใช้งานในรูปแบบ JSON บนเครื่องเซิร์ฟเวอร์ ออกแบบตามหลัก SOLID/GRASP ด้วย Port/Adapter Pattern
    \item \textbf{การตั้งค่าและความยืดหยุ่น:} ระบบรองรับการเพิ่มแหล่งข่าวใหม่ได้โดยไม่ต้องแก้ไขโครงสร้างหลัก ผ่านการลงทะเบียน Source ด้วย Decorator (\texttt{@register\_source})
    \item \textbf{ระบบแคชหลายระดับ (Multi-tier Cache):} ระบบมีกลไกจัดเก็บข้อมูลข่าวและเนื้อหาในหน่วยความจำ (\texttt{AsyncInMemoryCache}) ร่วมกับแคชแบบมีเงื่อนไข (HTTP 304 Conditional Cache) เพื่อลดการดึงข้อมูลซ้ำและเพิ่มความรวดเร็วในการแสดงผล
    \item \textbf{การรองรับ Progressive Web App (PWA):} รองรับการทำงานแบบเว็บแอปพลิเคชันก้าวหน้า พร้อม Service Worker สำหรับการแคชหน้าเว็บเพื่อความรวดเร็ว
\end{enumerate}

\subsection{ข้อจำกัดของโครงงาน}

\begin{enumerate}[label=\arabic*)]
    \item \textbf{ข้อจำกัดด้านทรัพยากรฮาร์ดแวร์และการติดตั้งใช้งาน:} ระบบถูกออกแบบและปรับแต่งให้รองรับภาระงานได้ตามขีดความสามารถของเครื่องแม่ข่าย AWS EC2 อินสแตนซ์ประเภท \texttt{c7i-flex.large} (2 vCPU, หน่วยความจำหลัก RAM ขนาด 4 GiB และ Swap Buffer ขนาด 4 GB) ซึ่งอาจมีข้อจำกัดในการรองรับผู้ใช้งานพร้อมกันจำนวนมหาศาล (Concurrent Users) หรือการรันเบราว์เซอร์ Playwright จำนวนหลายอินสแตนซ์พร้อมกัน
    \item \textbf{ข้อจำกัดของชุดข้อมูลฝึกสอนโมเดล (Training Dataset Constraints):} โมเดลการเรียนรู้ของเครื่องสำหรับการจำแนกหมวดหมู่ข่าวได้รับการฝึกสอนและประเมินผลบนชุดข้อมูลจำกัดจากคลังข้อมูล ThaiSum และ Prachathai-67k ซึ่งอาจมีความครอบคลุมของคำศัพท์เฉพาะกลุ่มหรือโครงสร้างภาษาตามลักษณะของแหล่งข่าวดังกล่าว
    \item \textbf{ความหลากหลายของการพาดหัวข่าวและแท็กหมวดหมู่ข้ามสำนักพิมพ์:} ข่าวเหตุการณ์เดียวกันจากต่างสำนักข่าวอาจมีการพาดหัวคนละมุมมอง หรือมีโครงสร้าง URL และ Tag หมวดหมู่ที่แตกต่างกัน ซึ่งระบบแก้ไขด้วยการใช้ฟีเจอร์แบบไฮบริด (หัวข้อ, เนื้อหาย่อ, คำสำคัญ, เอนทิตี) และการหาฉันทามติของหมวดหมู่ในระดับกลุ่มข่าว (Cluster-Level Category Consensus)
    \item ระบบรองรับการสรุปและแสดงผลข่าวในภาษาไทยและภาษาอังกฤษเท่านั้น ตามภาษาเดิมของเนื้อหาข่าวแต่ละฉบับ
    \item โครงสร้าง HTML หรือ RSS Feed ของเว็บไซต์แหล่งข่าวอาจเปลี่ยนแปลงได้ ทำให้ระบบต้องได้รับการปรับปรุงให้สอดคล้อง
    \item การสรุปข่าวต้องอาศัยการเชื่อมต่ออินเทอร์เน็ตและ API ของโมเดลภาษาขนาดใหญ่ (KKU Gen.AI) หากบริการขัดข้องหรือหมดโควตาการใช้งาน ระบบจะทำงานในโหมดข้อความสรุปตั้งต้นแทน
    \item รอบแรกของการดึงข่าวอาจใช้เวลา 2--5 นาที เนื่องจากต้องดึงเนื้อหาบทความของแต่ละข่าวเพื่อใช้ในการสรุป
\end{enumerate}

% -----------------------------------------------------------
\section{ประโยชน์ที่คาดว่าจะได้รับ}
% -----------------------------------------------------------

\begin{enumerate}[label=\arabic*)]
    \item ได้ระบบเว็บแอปพลิเคชันจำแนกและวิเคราะห์ข่าวสารอัตโนมัติแบบเรียลไทม์ ที่ช่วยให้ผู้ใช้ติดตามข่าวสารจากหลายแหล่งได้ในที่เดียวอย่างสะดวก รวดเร็ว และเหมาะสมกับความสนใจของแต่ละบุคคล
    \item ได้แนวทางในการประยุกต์ใช้เทคนิคเว็บสแครปปิง การอ่าน RSS Feed โมเดลภาษาขนาดใหญ่ อัลกอริทึมการกำหนดคุกกี้เพื่อการปรับแต่งเฉพาะบุคคล และอัลกอริทึมแนะนำข่าว เพื่อสร้างระบบวิเคราะห์ข้อมูลข่าวอัตโนมัติ ซึ่งสามารถนำไปพัฒนาต่อยอดกับข้อมูลประเภทอื่น
    \item ได้องค์ความรู้ด้านการประมวลผลภาษาธรรมชาติ การจำแนกข้อความด้วยกฎคำสำคัญร่วมกับโมเดลแมชชีนเลิร์นนิง เทคนิคการสรุปข้อความด้วยโมเดลภาษาขนาดใหญ่ และการพัฒนาระบบเว็บแบบเรียลไทม์
\end{enumerate}
